\section{总结与延伸}

本期从 Qwen3 的 Dense 和 MoE 两个变体入手，通过配置对比和参数量逐项计算，建立了对 Sparse MoE 架构的直观理解：

\begin{enumerate}
  \item MoE 将 FFN 替换为多个独立 Expert + Router，实现稀疏激活
  \item GQA 通过共享 KV 头压缩 KV Cache，是当前 LLM 的标配
  \item Qwen3 打破了传统的 $d_{\text{model}} = n_{\text{heads}} \times d_{\text{head}}$ 约束
  \item 30B 总参数中，推理时仅激活约 3.3B，推理效率接近同规模 Dense 模型
\end{enumerate}

\subsection{拓展阅读}

\begin{itemize}
  \item Sebastien Raschka 的博客：LLM 架构详解系列
  \item Qwen3 官方技术报告
  \item DeepSeek V2/V3 技术报告中的 MoE 设计
  \item 「Post NOS Chat GPT」系列中关于 LLaMA 2、GQA 的讲解
\end{itemize}

%% --- Fin del contenido --- %%

\end{document}
